\chapter{Introduction}\label{ch:introduction}

\section{Motivation}
Particle-laden channel flows connect local hydrodynamics with transport
outcomes. A particle can follow the carrier flow closely, cross a streamline
because of inertia, or approach a wall closely enough to change its final
destination. In a narrow passage, these outcomes depend on both the flow
field and the force closure. A small average velocity error can coexist
with a significant error in a particle's residence time or contact location.
Predictive modeling therefore requires more than visually plausible
streamlines.

Conventional computational methods provide established ways to discretize
the incompressible Navier--Stokes equations and integrate particle motion
\citep{ferziger2002,clift1978}. Physics-informed neural networks offer a
different representation: a differentiable network approximates the
solution, while differential equations and boundary conditions constrain
its parameters \citep{raissi2019,karniadakis2021}. This representation is
attractive when smooth field evaluation and derivatives are needed at
many moving particle positions. Its utility must nevertheless be established
by comparison with an independent reference and by checks that are sensitive
to the intended transport quantities.

I organize this study around a progressive increase in physical coupling
and geometric confinement. The first stage considers dilute particles whose
motion does not alter the carrier flow. The second introduces represented
particle populations and a spatially regularized reaction force in a straight
channel. The third defines a smooth constriction as the next geometry to be
studied. This progression isolates the consequences of adding two-way
coupling before introducing a varying channel height.

\section{Problem statement}
The central problem is to determine which parts of a physics-informed
fluid--particle model have sufficient numerical support for transport
prediction. There are three distinct questions. The first concerns the
carrier field: does the predicted velocity and pressure agree with a solution
of the same boundary-value problem? The second concerns the force transfer:
does the discrete coupling apply equal and opposite hydrodynamic forces
without introducing artificial interphase power? The third concerns
particle transport: are trajectories, contact events, and transit times
accurate for the geometry under consideration?

These questions are related but are not interchangeable. An algebraically
conservative transfer operator does not eliminate the residual error of a
trained fluid field. A low drag-network error does not correct an inaccurate
near-wall velocity. A successful straight-channel benchmark does not
establish performance at a constriction. I therefore treat each claim at
the level at which it has been evaluated.

\section{Aim and objectives}
The aim is to establish a physically consistent and numerically assessable
framework for physics-informed particle transport in channels. The study
has four objectives:
\begin{enumerate}
\item Formulate and evaluate a one-way coupled baseline with independent
carrier-flow and analytical drag references.
\item Derive a Gaussian parcel coupling that preserves force and interphase
work on its transfer quadrature.
\item Evaluate the components of the coupled formulation using independent
analytical and numerical verification problems.
\item Specify a constricted-channel extension whose geometry, boundary
conditions, and acceptance criteria can be tested without altering the
meaning of the earlier benchmarks.
\end{enumerate}

The corresponding research questions are: How accurately does the baseline
recover independently computed channel flow? Which conservation properties
follow from the Gaussian discretization, and which require additional
solution accuracy? What evidence is still needed before the formulation
can support predictions in a constricted channel?

\section{Scope and physical assumptions}
The completed channel models describe incompressible Newtonian flow and
spherical heavy particles moving in a plane. The baseline uses one-way
Stokes drag and a prescribed irreversible wall-sticking rule. The Gaussian
model represents spherical-particle populations through regularized force
distributions. The Gaussian cloud is a force-transfer representation; it
is not a resolved solid sphere and it is not an independently validated
description of a dense particle agglomerate.

The two-dimensional carrier does not resolve spanwise motion. No inference
about a finite-depth three-dimensional channel follows from an extruded
view. Particle rotation, lubrication, history force, and hydrodynamic
interactions of resolved finite-size bodies lie outside the accepted
baseline. The constriction is considered as a static wall geometry rather
than a dynamically growing deposit or a biological vessel.

\section{Study sequence and contributions}
\Cref{fig:route} summarizes the study sequence. The first contribution is a
baseline evaluation that combines analytical plane flow, an independently
assembled finite-element reference, a held-out drag test, and an explicit
distinction between field and trajectory accuracy. The second is a
consistent planar sphere-parcel formulation with shared Gaussian
interpolation and spreading. The third is a verification chain connecting
reduced relaxation problems, normalization identities, and an unsteady
incompressible-flow benchmark. The final contribution is a defined
constricted-channel study design that preserves the same requirements for
reference matching and conservation.

Completed findings are presented for the baseline carrier field and the
Gaussian model's verification problems. The coupled Gaussian channel
predictions have unresolved residual errors. The constricted-channel
formulation is prospective. This division determines the interpretation of
the conclusions and prevents a proposed geometry from being presented as
an executed calculation.

\ResultFigure[.92]{01_study_route}{Study progression. The first two panels
identify the model classes investigated; the third defines the proposed
constricted-channel extension. All three panels are geometric illustrations,
not simulation fields. Gaussian shading depicts kernel support rather than
a measured particle concentration.}{fig:route}

\section{Thesis organization}
\Cref{ch:literature} reviews the modeling choices and the requirements for
credible numerical evidence. \Cref{ch:framework} defines the common
nondimensional and evaluation framework. \Cref{ch:basic} develops the
baseline particle problem and reports its supported results.
\Cref{ch:gaussian} introduces the straight-channel Gaussian formulation and
its verification results. \Cref{ch:constriction} defines the constricted
channel and the proposed evaluation protocol. \Cref{ch:discussion} compares
the implications of the stages, and \cref{ch:conclusions} gives the
conclusions and priorities for further work. Supplementary derivations and
the computational record appear in the appendices.
